\section*{Capítulo \ref{ch:b2:equilibrium-shifts} --- Constantes de equilíbrio e deslocamentos}

\begin{solution}{exo:b2:equilibrium-shifts:1}
$K^\circ = \exp(32\,800/(8.314 \times 298.15)) = \num{5.6e5}$ e, para
$+\qty{32.8}{kJ/mol}$, seu inverso, \num{1.8e-6}. Um fator 100 exige
$\Delta(\Delta_r G^\circ) = -RT\ln 100 = \qty{-11.4}{kJ/mol}$.
\end{solution}

\begin{solution}{exo:b2:equilibrium-shifts:2}
$\ln K^\circ(400) = \ln(5.4 \times 10^5) + (91\,880/8.314)(1/400 -
1/298.15) = 13.20 - 9.44 = 3.76$, $K^\circ \approx 43$. As tabelas dão 35.6:
em \qty{100}{K}, a aproximação de $\Delta_r H^\circ$ constante já erra por
20\,\% (a reação fica mais \omterm{def:b2:reaction-enthalpy:exothermic}{exotérmica} à medida que $T$ aumenta).
\end{solution}

\begin{solution}{exo:b2:equilibrium-shifts:3}
(a) $n = 1$, $r = 0$, $\varphi = 2$: $v = 1$ (a pressão de vapor é função
de $T$). (b) $v = 3 - 1 + 2 - 3 = 1$. (c) $v = 3 - 1 + 2 - 1 = 3$. (d) Quatro
espécies, uma reação, duas fases (o carbono e o gás): $v = 4 - 1 + 2 - 2 =
3$; se o gás é formado apenas a partir de carbono e vapor d'água, $n(\ce{CO}) = n(\ce{H2})$
no gás, $s = 1$ e $v = 2$.
\end{solution}

\begin{solution}{exo:b2:equilibrium-shifts:4}
Aquecimento: sentido \omterm{def:b2:reaction-enthalpy:exothermic}{endotérmico}, para o \ce{NO2} (sentido direto). Compressão:
para menos moléculas de gás, o \ce{N2O4} (sentido inverso). Argônio a volume
constante: nenhuma mudança, as pressões parciais dos gases reagentes não variam.
\end{solution}

\begin{solution}{exo:b2:equilibrium-shifts:5}
A \qty{298}{K}: $\Delta_r G^\circ = 180.6 - 298.15 \times 0.0248 =
\qty{173.2}{kJ/mol}$, $K^\circ = \num{4.5e-31}$. A \qty{2000}{K}: $\Delta_r
G^\circ = 180.6 - 2000 \times 0.0248 = \qty{131.0}{kJ/mol}$, $K^\circ =
\num{3.8e-4}$. Com $\Delta\nu_{\text{gás}} = 0$, $x(\ce{NO})^2 = K^\circ x(\ce{N2})
x(\ce{O2})$: $x(\ce{NO}) = \sqrt{3.8 \times 10^{-4} \times 0.78 \times 0.21} =
\num{7.9e-3}$, quase um por cento. Os motores queimam a essas temperaturas;
à medida que os gases esfriam, $K^\circ$ cai enormemente, mas a decomposição
do \ce{NO} é muito lenta: o equilíbrio a quente fica congelado no
escapamento (até que um catalisador o elimine).
\end{solution}

\begin{solution}{exo:b2:equilibrium-shifts:6}
Com $\xi$ mol de éster e uma quantidade total constante, $Q =
\xi^2/[(1 - \xi)(n_B - \xi)]$. Para $n_B = 1.0$: $\xi^2/(1 - \xi)^2 = 4$, $\xi
= \qty{0.67}{mol}$. Para $n_B = 3.0$: $3\xi^2 - 16\xi + 12 = 0$, $\xi =
\qty{0.90}{mol}$. Retirar a água mantém $Q$ abaixo de $K^\circ$ qualquer que
seja $\xi$: a reação continua no sentido direto até o ácido ser consumido.
\end{solution}

\begin{solution}{exo:b2:equilibrium-shifts:7}
Sem argônio: quantidades $1 - \alpha$, $2\alpha$, total $1 + \alpha$;
$Q = 4\alpha^2/(1 - \alpha^2) = 0.148$, $\alpha = \sqrt{0.148/4.148} = 0.189$.
Com \qty{1}{mol} de argônio a \qty{1}{bar}: total $2 + \alpha$, $Q =
4\alpha^2/[(1 - \alpha)(2 + \alpha)] = 0.148$, logo $4.148\alpha^2 + 0.148\alpha
- 0.296 = 0$, $\alpha = 0.250$: diluir a pressão total constante abaixa as
pressões parciais, o que favorece o lado com mais moléculas de gás. A
volume constante, as pressões parciais dos gases reagentes não mudam:
$\alpha$ continua 0.189.
\end{solution}

\begin{solution}{exo:b2:equilibrium-shifts:8}
Somente o sólido: $n = 3$, $r = 1$, $s = 1$ ($p(\ce{NH3}) = p(\ce{HCl})$),
$\varphi = 2$: $v = 3 - 1 - 1 + 2 - 2 = 1$; escolhida $T$, a pressão fica
fixada (uma ``pressão de dissociação''). Com amônia adicionada, $s = 0$ e $v = 2$:
podem-se escolher $T$ e uma pressão; o produto $p(\ce{NH3})\,p(\ce{HCl})$ é
fixado por $T$.
\end{solution}

\begin{solution}{exo:b2:equilibrium-shifts:9}
$\Delta_r H^\circ = -R\ln(K_2/K_1)/(1/T_2 - 1/T_1) = -8.314 \times
\ln(0.0173)/(1/600 - 1/500) = \qty{-101}{kJ/mol}$, mais negativa que a
\qty{298}{K} (\qty{-91.9}{kJ/mol}), como previa a \omterm{thm:b2:reaction-enthalpy:kirchhoff}{lei de Kirchhoff}
($\Delta_r C_p^\circ < 0$).
\end{solution}

\begin{solution}{exo:b2:equilibrium-shifts:10}
$Q = \dfrac{n_{\ce{NH3}}^2\,n_{\text{tot}}^2}{n_{\ce{N2}}n_{\ce{H2}}^3}
\left(\dfrac{p^\circ}{p}\right)^2$. A $T$, $p$ constantes:
$\partial\ln Q/\partial n_{\ce{N2}} = -1/n_{\ce{N2}} + 2/n_{\text{tot}}$, positiva
quando $n_{\ce{N2}} > n_{\text{tot}}/2$, isto é, $x(\ce{N2}) > 1/2$. Então $Q$
supera $K^\circ$ e o equilíbrio se desloca no sentido inverso. A volume constante,
$\partial\ln Q/\partial n_{\ce{N2}} = -1/n_{\ce{N2}} < 0$: sempre no sentido direto.
\end{solution}

\begin{solution}{exo:b2:equilibrium-shifts:11}
$K^\circ = \exp(-1620/(8.314 \times 1100)) = 0.84$: $p(\ce{CO2}) =
\qty{0.84}{bar}$ no equilíbrio, isto é, $n = pV/RT = 0.84 \times 10^5 \times
0.0100/(8.314 \times 1100) = \qty{0.092}{mol}$ de gás. Com \qty{0.050}{mol}
de calcário, tudo se decompõe ($p = \qty{0.46}{bar} < K^\circ p^\circ$):
não há equilíbrio, há ruptura de equilíbrio. Com \qty{0.200}{mol}, equilíbrio:
\qty{0.092}{mol} de \ce{CO2} e de \ce{CaO}, restando \qty{0.108}{mol} de \ce{CaCO3}.
\end{solution}

\begin{solution}{exo:b2:equilibrium-shifts:12}
O leito 1 sai a cerca de \qty{890}{K} e 68\,\% de conversão, o leito 2 a cerca de
\qty{785}{K} e 92\,\%, o leito 3 a cerca de \qty{725}{K} e 98\,\%. Em um único
leito adiabático, o gás se aquece ao reagir até encontrar a curva de
equilíbrio, que, a essa temperatura, permite apenas cerca de 68\,\%. Manter o
catalisador frio o tempo todo não é uma opção, porque ele só funciona a
quente; resfriar entre os leitos mantém a velocidade e eleva o limite de
equilíbrio. O último leito é limitado pelo equilíbrio à temperatura mais baixa
em que o catalisador ainda funciona (as unidades então absorvem o trióxido e
passam o gás mais uma vez sobre um catalisador).
\end{solution}

\begin{solution}{pb:b2:equilibrium-shifts:1}
\textbf{1.} $\Delta_r H^\circ = \qty{-91.88}{kJ/mol}$, $\Delta_r S^\circ =
2(192.77) - 191.61 - 3(130.68) = \qty{-198.1}{J/(K.mol)}$, $\Delta_r G^\circ =
\qty{-32.8}{kJ/mol}$, $K^\circ = \num{5.6e5}$.
\textbf{2.} $K^\circ(700) = \exp(-54\,380/(8.314 \times 700)) = \num{8.8e-5}$;
$K^\circ(800) = \exp(-77\,320/(8.314 \times 800)) = \num{8.9e-6}$.
\textbf{3.} $\Delta_r H^\circ = -R\ln(K_{800}/K_{700})/(1/800 - 1/700)$,
isto é, \qty{-106}{kJ/mol}, mais negativa que a \qty{298}{K}.
\textbf{4.} $\ln K^\circ(723) = \ln(8.75 \times 10^{-5}) + 12\,777 \times
(1/723.15 - 1/700)$, with $12\,777 = 106\,230/8.314$: $K^\circ =
\num{4.9e-5}$.
\textbf{5.} $\Delta_r G^\circ(723) = -91.88 + 723.15 \times 0.1981 =
\qty{51.4}{kJ/mol}$, $K^\circ = \num{1.9e-4}$: quatro vezes grande demais, porque
$\Delta_r H^\circ$ e $\Delta_r S^\circ$ variam em \qty{425}{K}
($\Delta_r C_p^\circ \approx \qty{-44}{J/(K.mol)}$).
\textbf{6.} \ce{N2} $1 - \xi$, \ce{H2} $3 - 3\xi$, \ce{NH3} $2\xi$, total $4 -
2\xi$.
\textbf{7.} $x = 2\xi/(4 - 2\xi)$, logo $1 - x = (4 - 4\xi)/(4 - 2\xi)$; a
fração de nitrogênio $(1 - \xi)/(4 - 2\xi) = (1 - x)/4$, e a de hidrogênio
é três vezes maior.
\textbf{8.}
\[
  K^\circ = \frac{x^2}{\frac{1-x}{4}\left(\frac{3(1-x)}{4}\right)^3}
  \left(\frac{p^\circ}{p}\right)^2 = \frac{256}{27}\,\frac{x^2}{(1 - x)^4}
  \left(\frac{p^\circ}{p}\right)^2 ;
\]
tome a raiz quadrada.
\textbf{9.} $a = (\sqrt{27}/16) \times \sqrt{4.9 \times 10^{-5}} \times 1 =
\num{2.27e-3}$; $x/(1 - x)^2 = a$ dá $x = 0.0023$.
\textbf{10.} $a = 0.455$; $ax^2 - (2a + 1)x + a = 0$: $x = 0.25$.
\textbf{11.} $n = 3$, $r = 1$, $s = 1$ (razão 1 : 3 mantida pela reação),
$\varphi = 1$: $v = 2$. Escolhidos $T$ e $p$, nada mais fica livre.
\textbf{12.} Pressão: sentido direto ($\Delta\nu_{\text{gás}} = -2$). Temperatura:
sentido inverso (\omterm{def:b2:reaction-enthalpy:exothermic}{exotérmica}).
\textbf{13.} Um gás inerte a pressão total constante abaixa as pressões
parciais dos gases reagentes, como faria uma queda de pressão: sentido inverso.
\textbf{14.} $\partial\ln Q/\partial n_{\ce{N2}} = (-1/0.30 + 2)/n_{\text{tot}} < 0$:
$Q$ cai abaixo de $K^\circ$, sentido direto.
\textbf{15.} Não: a volume constante, $\partial\ln Q/\partial n_{\ce{N2}} =
-1/n_{\ce{N2}} < 0$ qualquer que seja a composição: sentido direto.
\textbf{16.} Quase sem amônia, $Q$ fica muito abaixo de $K^\circ$ na entrada:
o gás entra longe do equilíbrio e reage no sentido direto.
\textbf{17.} O gás passa tempo de menos sobre o catalisador para que o
equilíbrio seja atingido; um contato mais longo exigiria um conversor
maior e mais caro para um ganho pequeno.
\textbf{18.} $x = 2\xi/(4 - 2\xi) = 0.18$ dá $\xi = \qty{0.305}{mol}$ por
mol de \ce{N2} alimentado: \omterm{def:b2:equilibrium-shifts:single-pass}{conversão por passagem} do hidrogênio (e do nitrogênio)
$0.305$, 31\,\%.
\textbf{19.} Em regime estacionário, toda a alimentação nova acaba sendo convertida:
saem \qty{2.00}{mol} de amônia por unidade de tempo; o conversor precisa fazer
passar $1.00/
0.305 = \qty{3.3}{mol}$ de \ce{N2} por unidade de tempo, sendo o resto reciclado.
\textbf{20.} O argônio (e o metano) que entram com a alimentação nunca reagem e
não são removidos com a amônia líquida: sem \omterm{def:b2:equilibrium-shifts:single-pass}{purga}, eles se acumulariam no
circuito e abaixariam as pressões parciais dos reagentes.
\textbf{21.} A \qty{500}{K}, o catalisador de ferro é lento demais: o equilíbrio
seria favorável, mas nunca seria alcançado em um tempo industrial.
\textbf{22.} A \qty{723}{K} e \qty{200}{bar}, alimentação estequiométrica, gases
perfeitos: $\boldsymbol{x(\ce{NH3}) \approx 0.25}$.
\end{solution}
